\paragraph{Processing pipeline overview.}
\label{sec:pipeline_overview}

The proposed pipeline is an epoch-aware, threshold-based pipeline that transforms a continuous multichannel EEG recording into a non-overlapping epoch through a
fixed sequence of five operations.
The Figure~\ref{fig:pipeline_flow} summarises this
sequence as a flow chart, from the preprocessing, epoching, and the double thresholding which is the proposed method, through three detection stages (A, B, and C) which we will discussed in detail in section.




\begin{figure*}[!t]
    \centering
    \includegraphics[width=\textwidth]{flowchart}
    \caption{Epoch-aware double-thresholding pipeline for multichannel EEG blink-region detection. Continuous EEG signals from FP1, AF3, F3, FP2, AF4, and F4 are downsampled to 100 Hz, band-pass filtered from 1–20 Hz, and segmented into non-overlapping 30-s epochs after low-health epochs are removed. Stage A applies channel-specific peak-to-peak screening and a multichannel union rule to select suspicious epochs; E2, E5, and E7 form the candidate set in this example. Stage B pools and concatenates samples from these candidate epochs to estimate a robust sample-level threshold separately for each channel. Stage C applies the resulting channel-specific threshold to all valid epochs, including candidate and non-candidate epochs, and reports each detected blink region by epoch ID, channel, onset time, and duration. Orange denotes Stage A candidate epochs, while red denotes detected blink regions.}
    \label{fig:pipeline_flow}
\end{figure*}


\paragraph{Preprocessing}
First, the continuous recording is downsampled to a common sampling grid of \(100\)~Hz. Standardising the sampling rate makes the
amplitude thresholds and the minimum-duration gate comparable across recordings acquired on different hardware, and reduces the per-sample computational cost of the downstream stages without discarding any information needed to resolve a blink, whose
spectral content lies well below the new Nyquist frequency.

Second, a zero-phase band-pass filter from \(1.0\)~Hz  to \(20.0\)~Hz
is applied which is removes slow drifts and DC offsets at the low end and high-frequency muscle and line noise at the high end, isolating the low-frequency, high-amplitude deflection that characterises an eye blink.

Third, each filtered recording is split into non-overlapping, fixed-length epochs. Epoching is what makes the thresholds
\emph{adaptive}, rather than deriving a single level for the entire session, the pipeline treats each short epoch as a local statistical context.


%Finally, the detected regions are compared with the
%ground-truth annotations at the \emph{event} level using overlap-based matching, yielding
%precision, recall, and \(F_1\); these are macro-averaged across sessions and compared between
%configurations with Wilcoxon signed-rank tests, as detailed in
%Section~\ref{sec:eval}.
